% Shared configuration for Computational Mathematics notes (English).
\usepackage[T1]{fontenc}
\usepackage[utf8]{inputenc}
\usepackage[english]{babel}
\usepackage{lmodern}
\usepackage{microtype}

% Mathematics
\usepackage{amsmath,amssymb,amsthm}
\usepackage{mathtools}
\usepackage{bm}

% Algorithms: numbered lines and readable syntax
\usepackage[ruled,vlined,linesnumbered]{algorithm2e}
\SetKwInput{KwInput}{Input}
\SetKwInput{KwOutput}{Output}
\SetKwComment{tcp}{ // }{}
\DontPrintSemicolon

% Graphs and diagrams
\usepackage{tikz}
\usetikzlibrary{positioning,arrows.meta,trees,shapes.geometric,bending,calc,decorations.pathreplacing}

% Tables and layout
\usepackage{booktabs}
\usepackage{tabularx}
\usepackage{array}
\usepackage{float}
\usepackage{placeins}
\newcolumntype{Y}{>{\raggedright\arraybackslash}X}

% Source code
\usepackage{xcolor}
\usepackage{listings}
\definecolor{codebackground}{HTML}{F5F6F7}
\definecolor{codecomment}{HTML}{64748B}
\definecolor{codekeyword}{HTML}{1D4ED8}
\definecolor{codestring}{HTML}{0D9488}
\lstdefinestyle{code}{
  backgroundcolor=\color{codebackground},
  basicstyle=\ttfamily\small,
  commentstyle=\color{codecomment},
  keywordstyle=\color{codekeyword}\bfseries,
  stringstyle=\color{codestring},
  numbers=left,
  numberstyle=\scriptsize\color{codecomment},
  frame=single,
  breaklines=true,
  showstringspaces=false
}
\lstdefinestyle{pythoncode}{
  style=code,
  language=Python,
  morekeywords={as,self,True,False,None}
}

% Visual callout boxes
\usepackage[most]{tcolorbox}
\newtcolorbox{alertblock}[2][]{
  colback=red!5!white,
  colframe=red!75!black,
  fonttitle=\bfseries\small,
  title={#2},
  boxrule=0.8pt,
  arc=3pt,
  #1
}
\newtcolorbox{noteblock}[2][]{
  colback=blue!5!white,
  colframe=blue!75!black,
  fonttitle=\bfseries\small,
  title={#2},
  boxrule=0.8pt,
  arc=3pt,
  #1
}

% Formal environments in English (independent counters, numbered by section)
\theoremstyle{plain}
\newtheorem{theorem}{Theorem}[section]
\newtheorem{lemma}{Lemma}[section]
\newtheorem{property}{Property}[section]
\newtheorem{proposition}{Proposition}[section]
\newtheorem{corollary}{Corollary}[section]

\theoremstyle{definition}
\newtheorem{definition}{Definition}[section]
\newtheorem{example}{Example}[section]
\newtheorem{remark}{Remark}[section]

% Math macros
\newcommand{\R}{\mathbb{R}}
\newcommand{\C}{\mathbb{C}}
\newcommand{\N}{\mathbb{N}}
\newcommand{\norm}[1]{\left\|#1\right\|}
\newcommand{\abs}[1]{\left|#1\right|}
\newcommand{\inner}[2]{\left\langle #1, #2 \right\rangle}
\newcommand{\BigO}[1]{\mathcal{O}\!\left(#1\right)}
\DeclareMathOperator{\rank}{rank}
\DeclareMathOperator{\diag}{diag}
\DeclareMathOperator{\trace}{Tr}
\DeclareMathOperator{\cond}{cond}
\DeclareMathOperator{\spanop}{span}
\DeclareMathOperator*{\argmin}{arg\,min}
\DeclareMathOperator*{\argmax}{arg\,max}

% Hyperref
\usepackage[colorlinks=true,hypertexnames=false,linkcolor=blue!50!black,citecolor=blue!50!black,urlcolor=blue!50!black]{hyperref}
\hypersetup{
  pdftitle={Computational Mathematics for Learning and Data Analysis Notes (EN)},
  bookmarksopen=true
}
